\sscp{Bandicoots \it{(Peramelidae)}}{C-08-02}

\noindent{\wg\st{2pt}\begin{tabularx}{\linewidth}
{llll>{\raggedright\arraybackslash}Xl}
\lsptoprule
%Language	&	ISO	&	Term	&	Gloss	&	Source	\\	\midrule
%\endfirsthead
Language	&	ISO	&	Term	&	Gloss	&	Source	\\	\midrule
%\endhead
\ili{Matbat}		&	xmt	&	isɔ¹²k	&	bandicoot, \it{Echymipera kalubu}	&	\cite{Schapper2011}	\\
Ma{\Q}ya	&	slz	&	kalubu	&	bandicoot, \it{Echymipera kalubu}	&	\cite{Flannery1995}	\\
\ili{Biak}			&	bhw	&	karau (sop)	&	bandicoot, \it{Echymipera kalubu}	&	\cite{Schapper2011}	\\
\ili{Moor}			&	mhz	&	atíata	&	bandicoot	&	\cite{Schapper2011}	\\
\ili{Ansus}			&	and	&	pimunau	&	bandicoot	&	\cite{PriceDonohue2009}	\\
\ili{Ambai}\footnotemark			&	amk	&	kafa 	&	bandicoot (generic)	&	\cite{Price2002}	\\	\hhline{~}
%					&			&	kafa manduri;	&	species of bandicoot	&		\\	\hhline{~}
%					&			&	kafa maruba	&species of large bandicoot	&		\\
\lspbottomrule
\end{tabularx}}
\footnotetext{
		The following \ili{Ambai} terms for bandicoots are given in \cite{Price2002}:
		\it{kafa ne pari manduri} ‘Spiny Bandicoot,
		\it{Echymipera kalubu}’, \it{kafa manduri} ‘species of bandicoot’,
		\it{kafa maruba} ‘species of large bandicoot’.}

%\noindent{\wg\st{2pt}\begin{xltabular}{\linewidth}
%{lllll>{\raggedright\arraybackslash}X}
%\lsptoprule
%%Language	&	ISO	&	Term	&	Gloss	&	Etymon	&	Source	\\	\midrule
%%\endfirsthead
%Language	&	ISO	&	Term	&	Gloss	&	Etymon	&	Source	\\	\midrule
%\endhead
%
%\lspbottomrule
%\end{xltabular}}